
\section{Related Work}

\paragraph{Agentic Workflow}
Agentic workflow and autonomous agents ~\citep{zhuge2023mindstorms, hong2024data, jia2024mobile, wang2023voyager} represent two distinct paradigms of LLM application. The former completes tasks statically through predefined processes with multiple LLM invocations, while the latter solves problems dynamically through flexible autonomous decision-making. Compared to autonomous agents that require specific actions and decision patterns designed for the environment, agentic workflows can be constructed based on existing human domain experience and iterative refinement, offering higher potential for automated construction.

Agentic workflows can be broadly categorized into general and domain-specific types. General workflows emphasize universal problem-solving approaches, such as ~\citep{wei2022COT, wang2022COTSC, madaan2023self, wang2023multipersona}. Domain-specific workflows focus on building effective processes to solve domain-specific problems, such as code generation ~\citep{sirui2024meta, Tal2024Alpha, li2024debug}, data analysis ~\citep{yu2024hai, yi2024gen, bo2024nlsql, zhou2023llm}, mathematics ~\citep{zhong2024achieving, xu2024lemur}, question answering ~\citep{nori2023medprompt, zhou2024language}. Existing work has manually discovered numerous effective agentic workflows, but it's challenging to exhaust various tasks across different domains, further highlighting the importance of automated workflow generation and optimization. 

% Agentic workflow is a distinct paradigms with autonomous agents ~\citep{zhuge2023mindstorms, hong2024data, jia2024mobile, wang2023voyager}, with 前者通过predefined processes with 多次LLM调用静态的完成任务，后者通过flexible autonomous decision-making动态解决，这成为两种LLM应用的范式。相较于autonomous agents 需要针对环境设计具体的动作与决策模式, agentic workflow 的构建可以依赖人类已有的领域经验与迭代完成，具有更高的automated 构建的可能性。

% Agentic workflow 可以被简单分为General与Domain Specific两种类型，前者emphasize解决问题的通用思想，如~\citep{cot, stepback, xxxx}，后者关注于如何构建有效的流程以解决领域问题，如代码生成~\citep{sirui2024meta, Tal2024Alpha, li2024debug}, 数据分析~\citep{yu2024hai, yi2024gen, bo2024nlsql, alex2024automl}，数学~\citep{zhong2024achieving, }，Question Answering~\citep{}, add other field here. 现有的工作已经manually发掘了大量Effective  agentic workflow, 但难以穷尽不同领域的各种任务，进一步说明了对domain specific 任务 进行automatic workflow generate and optimization的重要性。

\paragraph{Automated Agentic Optimization}
Recent work aims to automate the design of agentic workflows, categorized into three types: automated prompt optimization, hyperparameter optimization, and automated workflow optimization. Prompt optimization ~\citep{chirs2024probreed, mert2024text, cheng2024opro, omar2024dspy} uses LLMs to optimize prompts within fixed workflows. Hyperparameter optimization \citep{saad2024archon} focuses on optimizing predefined parameters. While these approaches improve performance, they are limited in generalization to new tasks and often require moderate human effort for task-specific designs.

Automated workflow optimization~\citep{ze2024autoflow, wang2024symbo, zhuge2024gptswarm, hu2024automated} aims to optimize entire workflow structures, offering more potential for fully automated generation. Recent works explore diverse representations and methods. GPTSwarm~\citep{zhuge2024gptswarm} uses graph structures with reinforcement learning, but struggles to represent workflows with conditional states due to graph structure limitations. ADAS~\citep{hu2024automated} utilizes code structures to represent workflows and stores historical workflows in a linear list structure, aligning closely with our goals. However, it is constrained by the efficiency of its search algorithm as it relies on overly simplistic representations of experiences in the searching process, making it challenging to discover effective workflows.

\model also uses code to represent workflows, but goes further by providing a more fundamental structure called named node.  This structure encompasses various LLM invocation parameters, allowing for more detailed workflow representation. We also introduce operators that implement predefined node combination functions. Simultaneously, \model employs a specially designed MCTS algorithm for automated workflow optimization, leveraging the tree-structured experience and execution feedback to efficiently discover effective workflows.

% \model also uses code to represent workflows, but goes further than ADAS by providing a more fundamental structure called named node. This structure encompass various LLM invocation parameters, including model, prompt, temperature, and output format, allowing for more detailed workflow representation and expanding the space for automatic workflow optimization. 
% % 同时，我们引入了operator，能够通过与node一致的调用方式实现一系列预先定义好的node组合功能。
% Simultaneously, \model employs a specially designed MCTS algorithm for automatic workflow optimization, leveraging ree-structured experience and execution feedback to efficiently discover effective workflows.